\documentclass[10pt,letterpaper]{article}
\usepackage{cornell-notes}
\usepackage{mathtools}

\title{Chapter 21: Computational Geometry}
\author{}
\date{}

\setDocTitle{Chapter 21: Computational Geometry}
\setDocSubtitle{Numerical Methods}
\setDocVersion{v1.0}
\setDocStatus{Study Companion}
\setDocOwner{Numerical Methods Cornell Notes}
\setDocAuthor{}
\setDocDate{\today}
\setDocSummary{Cornell-format study and review notes for Chapter 21: Computational Geometry.}

\setCornellCollection{Numerical Methods Cornell Notes}
\setCornellUnitType{Chapter}
\setCornellUnitNumber{21}
\setCornellUnitTitle{Computational Geometry}
\setCornellCompanionLabel{Study and review companion}

\hypersetup{pdftitle={Chapter 21: Computational Geometry},pdfsubject={Cornell notes},pdfauthor={}}

\begin{document}
\maketitle

\section*{Chapter Overview}
Geometric predicates, proximity, and spatial structures define the heart of computational geometry. Reliable implementations check invariants before trusting a spatial algorithm.

\section*{Learning Objectives}
\begin{itemize}
\item Explain robust predicates, geometric representations, and spatial indexing.
\item Identify degeneracy, precision, and topology risks in a geometric algorithm.
\item Validate a result using an independent invariant or exact case.
\end{itemize}

\subsection*{Formula and notation anchor}
\[
\operatorname{orient}(a,b,c)=\det\!\begin{bmatrix}b_x-a_x&b_y-a_y\\c_x-a_x&c_y-a_y\end{bmatrix}
\]

\begin{exambox}{Numerical warning}
Near-collinearity, inconsistent boundary conventions, and approximate predicates can corrupt downstream topology.
\end{exambox}

\section{21.0 Introduction}
\begin{cornell}
\noterow{What is the central idea?}{Computational geometry turns spatial relationships into predicates and constructions that must remain reliable near degeneracy.}
\noterow{How should it be used?}{Define coordinate scale and boundary conventions, separate predicates from constructions, and validate topology and geometric invariants.}
\noterow{Cost and reliability}{Analyze arithmetic work, storage, and search pruning separately. Use robust predicates or adaptive precision when a sign decision is near zero.}
\noterow{Implementation checklist}{\begin{itemize}
\item Validate dimensions, domains, ordering, and structural assumptions.
\item Guard degeneracy and nonfinite coordinates.
\item Test exact, boundary, and difficult valid cases.
\end{itemize}}
\end{cornell}

\begin{summarybox}{Section summary}
Select a geometric representation and algorithm that match the data, then accept the result only after a residual, invariant, or topology check that is independent of the implementation path.
\end{summarybox}

\end{document}
